\documentclass[10pt,a4paper]{amsart}
\usepackage[margin=19mm]{geometry}
\usepackage{amsmath,amssymb,mathtools}
\usepackage[scr=rsfs,cal=euler]{mathalfa}
\usepackage{xfrac}
\usepackage[unicode,hidelinks,bookmarks=false]{hyperref}
\newcommand{\ie}{i.e.\ }
\newcommand{\cf}{cf.\ }
\newcommand{\resp}{respectively\ }
\newcommand{\Subsection}{\subsection}
\providecommand{\nobreakdash}{\nobreak\hbox{-}}
\setlength{\emergencystretch}{2em}
\multlinegap0pt
\usepackage{amssymb}
\usepackage{mdframed}
\usepackage{bm}
\usepackage{xcolor}
\newtheorem{thm}{Theorem}
\newtheorem{lemma}{Lemma}
\newcommand{\calH}{\mathcal{H}}
\title{Arithmetic geometry of quantum connections on Calabi--Yau $3$-folds}
\author{Shaoyun Bai, Jae Hee Lee, Daniel Pomerleano}
\date{Source: arXiv:2601.01654v2 (CC BY 4.0)}
\begin{document}
\maketitle

\noindent\emph{Source and licence.} Arithmetic geometry of quantum connections on Calabi--Yau $3$-folds. Version arXiv:2601.01654v2 (CC BY 4.0), distributed by arXiv under the Creative Commons Attribution 4.0 International licence, http://creativecommons.org/licenses/by/4.0/, as recorded in the arXiv licence metadata for this exact version and shown on the abstract page https://arxiv.org/abs/2601.01654v2. Full licence notice: \texttt{licenses/arxiv-2601.01654-CC-BY-4.0.txt}.\par\smallskip

\noindent\textit{Editorial scope.} Counted unit: the article's thm thm:integrality (source0.tex lines 297--297). The article's complete original proof of it is delivered with it (source0.tex lines 1075--1089, one \texttt{proof} environment of 15 lines, which the article places away from the statement). No mathematics is written, repaired or re-derived: the statement and the proof are the article's own lines, in the article's own notation, and no retained line is substituted or re-flowed.  Retained uncounted as the source-local context the proof needs: lines 292--293; lines 942--942; lines 946--946; lines 1050--1060; lines 2390--2396.  Nothing was omitted from any retained span, so the package carries no omission record.  No retained line contains a citation, so the wrapper carries no bibliography and no bibliography entry.  No retained line contains a figure environment, an image-inclusion command or drawing code, so no image is delivered and the package declares no asset dependency.  Editorial formatting only, in addition: an AMS \texttt{amsart} preamble that restates the article's own declarations and macros used by the retained lines, namely the environments \texttt{thm}, \texttt{lemma} on separate counters, together with the standard packages those lines need.  The retained environments print in this document as Theorem 1, Lemma 1 in order of appearance, whereas the article prints the same items with the numbering of its own full document.  The complete native source of the article as distributed in the arXiv e-print for this exact version is bundled unchanged as \texttt{sources/192097/source0.tex}.
We use $\Gamma_p(TX)$ to introduce a grading preserving $\mathbb{Z}_p$-linear map on cohomologies: \begin{align} \label{eq:b-frobenius} 
\Phi_0: H^*(X;\mathbb{Z}_p)[u,u^{-1}] \to H^*(X;\mathbb{Q}_p)[u,u^{-1}], \quad \Phi_0(u^kx) = p^{3+k} u^k\cdot(\Gamma_p(TM) \smile x). \end{align} 

\begin{align} \label{eq:differentialequations} q_k\partial_{q_k}(\Phi(q_1, \cdots, q_r)) + M_k(q_1, \cdots, q_r)\Phi(q_1, \cdots, q_r)-p\Phi(q_1, \cdots, q_r)M_k(q_1^p, \cdots, q_r^p)=0. \end{align}

\begin{lemma} \label{lem:initialtermdeterminesFrob} Fix a matrix $\Phi_0 \in \operatorname{Mat}_{s\times s}(\mathbb{Q}_p(q_1,\cdots,q_{r-1})).$ There is at most one solution $\Phi(q)$ to  \eqref{eq:differentialequations} with coefficients in $L$ whose leading order term is $\Phi_0.$ \end{lemma}

\begin{lemma} \label{lem:integralityusingKSV} 

Fix $\zeta \in \mathbb{Q}_p$ with p-adic valuation $\operatorname{val}_p(\zeta)>-3.$ Fix as well the preferred homogeneous basis for $\calH^0$ given by $$1 \in H^0(X), u^{-1} e_i \in H^2(X), u^{-2}L_i \in H^4(X), u^{-3}[\mathrm{pt}] \in H^6(X).$$  Define a map $\Phi$ on this basis by the formula: 
\begin{align}
\Phi (1) &= p^3 1 +p^3 u^{-2} \sum_j\delta_j(\eta_\zeta (q)) L_j -2p^3 u^{-3} \eta_\zeta (q) [\mathrm{pt}], \\
\Phi (u^{-1}e_i) &= p^2 (u^{-1}e_i) + p^2 u^{-2}  \sum_j \delta_{i}\delta_{j}(\eta_\zeta(q))L_j -p^2 u^{-3}\delta_i(\eta_\zeta (q)) [\mathrm{pt}],\nonumber \\
\Phi (u^{-2}L_i) &= p \cdot u^{-2}L_i, \\
\Phi (u^{-3}[\mathrm{pt}])&=u^{-3}[\mathrm{pt}]. \nonumber  
\end{align}
The above map defines a horizontal Frobenius intertwiner: 
\begin{align}\label{eq:PhiH0} \Phi :  F^*\calH^0 \to \calH^0.\end{align} \end{lemma} 

\begin{lemma}\label{lem:pintegrality}
Let $p$ be an odd prime.
\begin{enumerate}
    \item[(a)] If $p \ge 5$, then $\zeta_p(3) \in \mathbb{Z}_p$.
    \item[(b)] If $p = 3$, then $v_3(\zeta_3(3)) = -1$ (in particular, $\zeta_3(3) \notin \mathbb{Z}_3$).
\end{enumerate}
\end{lemma}

\begin{thm}\label{thm:integrality} Let $p$ be a prime for which $H^{even}(X;\mathbb{Z}_p)$ contains no torsion. Then there is a unique Frobenius map: \begin{align}\label{eq:Amodelcrystalline} \Phi: F^*QH^*(X; \Lambda_{\mathbb{Z}_p})[u^{\pm}] \to QH^*(X; \Lambda_{\mathbb{Z}_p})[u^{\pm},1/p] \end{align} such that: \begin{enumerate} \item Let $\Lambda_{+} \subset \Lambda_{\mathbb{Z}_{p}}$ denote the ideal of elements with strictly positive $\omega$-valuation. The ``leading order term" (with respect to the $\omega$-valuation) of $\Phi$ is given by \eqref{eq:b-frobenius}: \begin{align} \Phi = \Phi_0  \quad \operatorname{mod}\ \Lambda_{+}.\end{align}  \item Let $\mathcal{H}^{2k}$ denote the subspace of $QH^*(X; \Lambda_{\mathbb{Z}_p})[u^{\pm}]$ of degree $2k$. The map $\Phi$ satisfies:   \begin{align} \label{eq:keyintegralitycondition} \Phi(F^*\mathcal{H}^{2k}) \subset p^k\mathcal{H}^{2k},\quad k \in \mathbb{Z}.\end{align} \end{enumerate} \end{thm}

\begin{proof}[Proof of Theorem \ref{thm:integrality}]
     For $\Phi$ as in Lemma \ref{lem:integralityusingKSV}, if we set $q=0$, we see that its reduction modulo the ideal $\Lambda_{+}$, $\Phi_0 $, satisfies
    \begin{equation}
        \begin{aligned}
            \Phi_0(1) &= p^31 -2p^3u^{-3}\zeta[pt], \\
            \Phi_0(u^{-1}[e_i]) &= p^2(u^{-1}[e_i]), \\
            \Phi_0(u^{-2}[L_i]) &= pu^{-2}[L_i], \\
            \Phi_0(u^{-3}[pt]) &= u^{-3}[pt].
        \end{aligned}
    \end{equation}
    We set $\zeta = \frac12 \zeta_p(3)\chi(X)$. By Lemma \ref{lem:pintegrality}, this satisfies the necessary condition on $\operatorname{val}_p(\zeta)$. Note that multiplication by $u^k$ gives identifications: \begin{align}\calH^{0} \cong \calH^{2k},\quad F^*\calH^{0} \cong F^*\calH^{2k}.
    \end{align} 
   We extend $\Phi$ to $\calH^{2k}$ by requiring that \begin{align} \Phi(u^k\cdot x)=p^ku^k \cdot \Phi(x),\quad x \in \calH^{0}. \end{align}
   This is again a horizontal Frobenius (on each piece $\calH^{2k}$ being just a rescaling of the Frobenius on $\calH^{0}$) and now has leading order term  \eqref{eq:b-frobenius}. By construction, it also satisfies: \begin{align} \Phi(F^*\calH^{2k}) \subset p^k \calH^{2k} \end{align} as required. Lastly, uniqueness is guaranteed by Lemma \ref{lem:initialtermdeterminesFrob}.
\end{proof}

\end{document}
